\chapter{Testing and Verification}

\section{Testing Methodology}

The processor verification follows a systematic approach:

\begin{enumerate}
    \item \textbf{Unit Testing}: Test individual modules (ALU, register file, memory)
    \item \textbf{Integration Testing}: Test complete processor with instruction sequences
    \item \textbf{ISA Testing}: Verify all instruction types with comprehensive test suite
    \item \textbf{System Testing}: Run TA-provided testbench for final verification
\end{enumerate}

\section{Testbench Structure}

The testbench consists of three main components:

\begin{itemize}
    \item \textbf{Driver}: Generates clock, reset, and input stimuli
    \item \textbf{DUT (Device Under Test)}: The processor module
    \item \textbf{Scoreboard}: Monitors outputs and verifies correctness
\end{itemize}

\section{Test Programs}

\subsection{ISA Test Suite}

The TA-provided test suite (\texttt{isa.mem}) includes tests for all instruction types:

\begin{itemize}
    \item Arithmetic instructions: ADD, ADDI, SUB
    \item Logical instructions: AND, ANDI, OR, ORI, XOR, XORI
    \item Comparison instructions: SLT, SLTI, SLTU, SLTIU
    \item Shift instructions: SLL, SLLI, SRL, SRLI, SRA, SRAI
    \item Load instructions: LW, LH, LHU, LB, LBU
    \item Store instructions: SW, SH, SB
    \item Branch instructions: BEQ, BNE, BLT, BGE, BLTU, BGEU
    \item Jump instructions: JAL, JALR
    \item Upper immediate instructions: LUI, AUIPC
\end{itemize}

\subsection{Expected Test Results}

A correctly functioning processor should produce the following output:

\begin{verbatim}
SINGLE CYCLE TESTS

add......PASS
addi.....PASS
sub......PASS
and......PASS
andi.....PASS
or.......PASS
ori......PASS
xor......PASS
xori.....PASS
slt......PASS
slti.....PASS
sltu.....PASS
sltiu....PASS
sll......PASS
slli.....PASS
srl......PASS
srli.....PASS
sra......PASS
srai.....PASS
lw.......PASS
lh.......PASS
lhu......PASS
lb.......PASS
lbu......PASS
sw.......PASS
sh.......PASS
sb.......PASS
auipc....PASS
lui......PASS
beq......PASS
bne......PASS
blt......PASS
bltu.....PASS
bge......PASS
bgeu.....PASS
jal......PASS
jalr.....PASS
malgn....ERROR (expected, misaligned not required)
iosw.....PASS

END of ISA test
\end{verbatim}

\section{Verification Strategy}

\subsection{Signal Monitoring}

The testbench monitors key signals:

\begin{itemize}
    \item \texttt{o\_pc\_debug}: Tracks program counter to verify instruction flow
    \item \texttt{o\_io\_ledr}: Used for test result output (scoreboard checks at specific PC values)
    \item \texttt{o\_insn\_vld}: Ensures processor is executing valid instructions
\end{itemize}

\subsection{Test Execution Flow}

\begin{enumerate}
    \item Testbench loads \texttt{isa.mem} into instruction memory (or \texttt{isa\_1b.hex} as fallback)
    \item Reset is applied, processor starts at PC = 0x0000\_0000
    \item Processor executes instructions sequentially
    \item Scoreboard monitors PC and LEDR outputs
    \item When PC reaches 0x18, LEDR[7:0] contains test name (ASCII)
    \item When PC reaches 0x1C, test completion is signaled
\end{enumerate}

\section{Simulation Tools}

The processor can be simulated using:

\begin{itemize}
    \item \textbf{Verilator}: Fast compilation and simulation
    \item \textbf{Xcelium}: Cadence simulator with GUI support
    \item \textbf{ModelSim/QuestaSim}: Mentor Graphics simulator
    \item \textbf{Vivado Simulator}: Xilinx simulator
\end{itemize}

\subsection{Verilator Workflow}

\begin{verbatim}
cd 10_sim
make              # Compile and run
make wave         # Open GTKWave for waveform viewing
\end{verbatim}

\subsection{Xcelium Workflow}

\begin{verbatim}
cd 11_xm
module load xcelium
make              # Run simulation
make gui          # Open SimVision GUI
\end{verbatim}

\section{Debugging Methodology}

When tests fail, debugging follows these steps:

\begin{enumerate}
    \item \textbf{Waveform Analysis}: Check signal values at failing instruction
    \item \textbf{PC Tracking}: Verify instruction fetch at correct address
    \item \textbf{Register Check}: Verify register file reads and writes
    \item \textbf{ALU Verification}: Check ALU inputs and outputs
    \item \textbf{Memory Access}: Verify memory reads and writes
    \item \textbf{Control Signals}: Verify control unit outputs
\end{enumerate}

\section{Common Issues and Solutions}

\subsection{Issue 1: Processor Hanging}

\textbf{Symptoms}: No test output, simulation runs indefinitely.

\textbf{Causes}:
\begin{itemize}
    \item Branch instruction logic incorrect
    \item PC update logic missing
    \item Infinite loop in test program
\end{itemize}

\textbf{Solution}: Check branch condition evaluation and PC update logic.

\subsection{Issue 2: Wrong Test Results}

\textbf{Symptoms}: Tests show ERROR instead of PASS.

\textbf{Causes}:
\begin{itemize}
    \item ALU operation incorrect
    \item Immediate generation wrong
    \item Memory access address calculation error
\end{itemize}

\textbf{Solution}: Verify instruction decoding and execution logic.

\subsection{Issue 3: Memory Access Errors}

\textbf{Symptoms}: Load/store operations fail.

\textbf{Causes}:
\begin{itemize}
    \item Wrong byte ordering (big-endian vs little-endian)
    \item Sign extension incorrect
    \item Address alignment issues
\end{itemize}

\textbf{Solution}: Verify byte ordering and sign extension logic.

\section{Code Coverage}

While formal code coverage analysis was not performed, the test suite ensures:

\begin{itemize}
    \item All instruction types are tested
    \item Both signed and unsigned operations are verified
    \item All branch conditions are exercised
    \item Memory-mapped I/O is tested
    \item Edge cases (register x0, misaligned addresses) are handled
\end{itemize}

\section{Performance Metrics}

The single-cycle processor has the following characteristics:

\begin{itemize}
    \item \textbf{Clock Frequency}: Determined by longest combinational path
    \item \textbf{Cycles per Instruction}: 1 (by design)
    \item \textbf{Critical Path}: Through ALU, memory access, or register file
    \item \textbf{Area}: Function of memory size and register file
\end{itemize}

These metrics are not critical for single-cycle implementation but will be important for pipelined versions in future milestones.

